% !TeX program = pdflatex
\documentclass[11pt]{article}

\usepackage[utf8]{inputenc}
\usepackage[T1]{fontenc}
\usepackage{amsmath,amssymb,amsthm}
\usepackage[margin=1in]{geometry}
\usepackage{booktabs}
\usepackage{graphicx}
\usepackage{svg}
\usepackage{hyperref}
\usepackage{microtype}

\hypersetup{colorlinks=true, linkcolor=blue, citecolor=blue, urlcolor=blue}

\newtheorem{theorem}{Theorem}[section]
\newtheorem{lemma}[theorem]{Lemma}
\newtheorem{proposition}[theorem]{Proposition}
\newtheorem{corollary}[theorem]{Corollary}
\theoremstyle{definition}
\newtheorem{definition}[theorem]{Definition}
\newtheorem{example}[theorem]{Example}
\newtheorem{problem}[theorem]{Problem}
\theoremstyle{remark}
\newtheorem{remark}[theorem]{Remark}
\newtheorem{conjecture}[theorem]{Conjecture}

\newcommand{\Cyc}{\mathrm{Cyc}}
\newcommand{\Per}{\mathrm{Per}}
\newcommand{\dens}{\mathrm{dens}}
\DeclareMathOperator{\ord}{ord}
\DeclareMathOperator{\rad}{rad}
\DeclareMathOperator{\lcmop}{lcm}

\title{\textbf{Iterated Digit-Sum Power Maps:\\ A Lower Bound Theorem, an Exact Basin Density Law,\\ and the Oscillation of Intra-Signature Basin Splits}}
\author{Alex Martins\\ \small Independent Research}
\date{August 2026 (v2.1.0; Tier B empirical patch, September 2026; see Appendix~B)}

\begin{document}
\maketitle

\begin{center}
\small \textbf{MSC 2020:} 37P35, 11A63, 11K06, 11A25, 05C20\\
\small \textbf{Keywords:} discrete dynamical systems, digit-sum functions, generalized happy numbers, modular arithmetic, attractor basins, natural density, equidistribution, functional graphs, digital functions.\\
\small \textbf{Zenodo:} \url{https://zenodo.org/records/22181953} \quad \textbf{DOI:} \url{https://doi.org/10.5281/zenodo.22181953}
\end{center}

\begin{abstract}
I study the family of discrete dynamical systems $f_{k,b}(n) = S_b(n^k)$, where $S_b$ is the base-$b$ digit-sum function and $k \ge 1$ a fixed exponent. The map is related to, but distinct from, the generalized happy functions $S_{e,b}(n) = \sum_i d_i^{\,e}$ of Grundman and Teeple; the distinction, and the modular structure that follows from it, is made precise in \S1. Every orbit is eventually periodic, and the system is organized by an algebraic skeleton: the functional graph of the modular power map $\varphi_{k,b-1}(x) = x^k \bmod (b-1)$.

I prove three structural results. \textbf{(I) Lower Bound Theorem.} The number of attractors satisfies $|C(k,b)| \ge \Cyc(\varphi_{k,b-1})$, the number of cycles of the modular map. \textbf{(II) Basin Density Law.} For each modular cycle $\gamma_i$ with residue basin $R_i$, the aggregate basin $\mathcal{N}_i$ of all attractors carrying the residue signature $\gamma_i$ has a natural density, equal exactly to the modular weight $p_i = |R_i|/(b-1)$; over a window $[1,M]$ the empirical proportion deviates by at most $\min(|R_i|,\,b-2)/M$, and the underlying finite-window statement is an exact identity between integers. \textbf{(III)} Both counting functions of $\varphi_{k,m}$ admit closed forms: the periodic-point count is the classical $\prod_{p^e \| m}\big(1 + \kappa_k(\varphi(p^e))\big)$, with $\kappa_k(N)$ the largest divisor of $N$ coprime to $k$, while the cycle count (the quantity required by the lower bound) is obtained by combining local cycle types through the Chinese remainder theorem. The distinction matters and was missed in an earlier draft: the periodic-point count depends on $k$ only through $\rad(k)$, whereas the cycle count does not.

Results (I) and (II) are elementary consequences of the digit-sum congruence (``casting out nines'') and of residue equidistribution. Novelty is claimed for the organization of these consequences and for isolating the remaining problem: how the mass $p_i$ splits among several physical attractors that share a signature. That partition is not elementary. I resolve it empirically and mechanistically. No individual basin density appears to exist; along fixed digit-length the masses oscillate quasi-periodically, in antiphase between competing attractors, while always summing to $p_i$. A parameter-free model, the Gaussian law of $S_b(n^k)$ (local limit theorem) on the image lattice $v \equiv r^k \pmod{b-1}$, convolved with the exact integer-to-attractor labelling, reproduces the oscillation to within Monte-Carlo noise (mean per-band bridge error $|\delta_j-F_j| \approx 0.0017$ on the $(3,10)$ pilot over $D=4\ldots 90$; \texttt{scripts/bridge\_check.py}). The oscillation is an instance of the classical log-periodic fluctuation of digital functions (Delange; Drmota--Grabner).

All claims are verified computationally: an exhaustive sweep of $19{,}500$ parameter pairs $(k,b)$ with $k \le 500$, $b \le 40$, with zero lower-bound violations and with the exact integer form of the aggregate law holding without exception, together with an independent re-implementation reported in Appendix~A.

\smallskip
\noindent\emph{Erratum.} This version corrects five errors in the July 2026 draft: a false statement about the modular cycle count, a misclassification of which signatures oscillate, an overstated claim of sharpness for the error bound, a confounded statistic, and a wrong congruence lattice in the Gaussian split model. Each is documented in Appendix~B, with the reasoning that produced it, since the failure modes are instructive.
\end{abstract}

\section{Introduction}

\subsection{The map, and how it differs from happy functions}

Fix a base $b \ge 2$ and set $m = b-1$. The \textbf{generalized happy function} studied by Grundman and Teeple \cite{GT01,GT07} and surveyed in \cite{survey} is
\[
S_{e,b}(n) \;=\; \sum_i d_i^{\,e}, \qquad n = \sum_i d_i\, b^i,
\]
the sum of the $e$-th powers of the base-$b$ \emph{digits} of $n$. The object of the present work is the different map
\[
f_{k,b}(n) \;=\; S_b\big(n^k\big),
\]
the base-$b$ \textbf{digit sum of $n^k$}. These are distinct dynamical systems. For $(k,b) = (2,10)$, the happy map sends $13 \mapsto 1^2+3^2 = 10 \mapsto 1$ (so $13$ is happy), whereas $f_{2,10}(13) = S_{10}(169) = 16$ and $f_{2,10}(16) = S_{10}(256) = 13$: a $2$-cycle. The classic happy-number iteration and $f_{2,10}$ are not the same system.

What the two families share is the arithmetic that drives everything below: the classical congruence
\begin{equation}\label{eq:co9}
S_b(N) \equiv N \pmod{m}, \qquad m = b-1,
\end{equation}
(``casting out nines''), valid for every $N$, since $b \equiv 1 \pmod{m}$ implies $d_i b^i \equiv d_i \pmod m$ digit by digit. For $f_{k,b}$ this yields the identity
\begin{equation}\label{eq:inv}
f_{k,b}(n) \;=\; S_b(n^k) \;\equiv\; n^k \pmod{m},
\end{equation}
so reducing the dynamics mod $m$ gives exactly the modular power map $\varphi_{k,m}(x) = x^k \bmod m$. (For the happy map $S_{e,b}$ the mod-$m$ reduction is $\sum_i d_i^{\,e}$, which is not a function of $n \bmod m$ alone; this is why $f_{k,b}$ has a simpler modular skeleton than the happy family.) The invariance \eqref{eq:inv} is elementary and classical. Novelty is not claimed for it, only for the structural and distributional consequences assembled below, and for isolating the split problem of \S\ref{sec:split}.

\subsection{The questions}

A systematic treatment of the family $f_{k,b}$ parametrized by exponent $k$ and base $b$ appears to be absent from the literature. The questions addressed here are:
\begin{enumerate}
\item \textbf{Count.} How many distinct attractors does $f_{k,b}$ possess, and can the count be bounded from the algebraic parameters alone? (Answer: the Lower Bound Theorem, \S\ref{sec:lb}.)
\item \textbf{Distribution.} How much mass, in the sense of natural density, does each attractor's basin carry? (Answer, at the aggregate level: the Basin Density Law, \S\ref{sec:bdl}.)
\item \textbf{Split.} When several attractors share a residue signature, how does the (exactly determined) aggregate mass divide among them? (Answer, empirical and mechanistic: the split does not settle; it oscillates; \S\ref{sec:split}.)
\end{enumerate}

The classic instance $(k,b) = (2,10)$ already displays the distributional phenomenon. The three attractors $\{1\}$, $\{9\}$, $\{13,16\}$ capture positive integers in proportions $22.2\%,\,33.3\%,\,44.4\%$, precisely $2/9,\,3/9,\,4/9$, the sizes of the residue classes modulo $9$ feeding each attractor. Section~\ref{sec:bdl} shows this is forced by \eqref{eq:inv} together with equidistribution, hence a law for aggregate mass, not a coincidence. Section~\ref{sec:split} shows that the freedom the law leaves untouched, the intra-signature split, is where the remaining difficulty lies.

The following separation is used throughout:
\begin{itemize}
\item \textbf{Determined by modular arithmetic (exactly, elementarily):} a lower bound on the number of attractors, and the total basin mass per residue signature.
\item \textbf{Not determined by modular arithmetic:} the excess attractor count $\Delta(k,b)$, and whether each signature's mass splits into convergent individual densities.
\end{itemize}

\subsection{Scope and contribution}

The theorems on aggregate mass and on cycle counts are short corollaries of the digit-sum congruence and of residue equidistribution. A reader fluent in casting-out-nines will find the proofs brief. The article records four things.

Theorem~\ref{thm:lb} gives the lower bound $|C(k,b)| \ge \Cyc(\varphi_{k,b-1})$. Theorem~\ref{thm:bdl} gives the aggregate basin density of each residue signature exactly, $p_i = |R_i|/(b-1)$, with the finite-window bound $\min(|R_i|,\,b-2)/M$ and with an exact integer identity underneath (Proposition~\ref{prop:partition}). Proposition~\ref{prop:cyc} gives a closed form for the cycle count of the modular map, which is the quantity Theorem~\ref{thm:lb} requires and which is not a function of $\rad(k)$. Remark~\ref{rem:result} records that, when several attractors share a signature, the individual basin densities do not appear to exist: along fixed digit-length the masses oscillate quasi-periodically, in antiphase, always summing to $p_i$.

The aggregate law turns patterns such as $22/33/44$ into a theorem rather than folklore, and it isolates the remaining object, the intra-signature split. That split is an instance of the log-periodic fluctuation of digital functions (Delange \cite{delange}, Drmota--Grabner \cite{drmota-grabner}). The periodic-point formula of \S\ref{sec:count} is classical power-digraph theory (Somer--K\v{r}\'i\v{z}ek \cite{somer-krizek}, Chou--Shparlinski \cite{chou-shparlinski}). The cycle-count formula of the same section is elementary; I have not found it stated for this purpose.

Both exact theorems were checked without exception on an exhaustive grid of $19{,}500$ pairs $(k,b)$ with $k \le 500$ and $b \le 40$ (Appendix~A). This version corrects five errors of the July 2026 draft (Appendix~B).

\section{Definitions and Notation}

Fix $b \ge 2$ and set $m = b-1$.

\begin{definition}[Digit sum]
For $n \ge 1$, $S_b(n)$ is the sum of the digits of $n$ in the standard base-$b$ representation.
\end{definition}

\begin{definition}[The map]
The \textbf{iterated digit-sum power map} is $f_{k,b}\colon \mathbb{Z}^+ \to \mathbb{Z}^+$, $f_{k,b}(n) = S_b(n^k)$.
\end{definition}

\begin{definition}[Attractors]
The \textbf{orbit} of $n$ is $(f_{k,b}^{\,t}(n))_{t \ge 0}$. An \textbf{attractor} is a minimal periodic orbit: a finite set $A = \{a_1,\dots,a_L\}$ with $f(a_i) = a_{i+1}$ (indices mod $L$), no proper subset periodic. $C(k,b)$ denotes the set of all attractors of $f_{k,b}$.
\end{definition}

\begin{definition}[Modular power map]
$\varphi_{k,m}(x) = x^k \bmod m$ on $\mathbb{Z}/m\mathbb{Z}$; $G(\varphi_{k,m})$ is its functional graph (every vertex has out-degree $1$, so each component is a $\rho$-shape: a cycle with in-trees). Its cycles are $\gamma_1,\dots,\gamma_c$, $c = \Cyc(\varphi_{k,m})$.
\end{definition}

\begin{definition}[Modular basin and weight]
For each cycle $\gamma_i$,
\[
R_i = \{\, r \in \mathbb{Z}/m\mathbb{Z} : \varphi_{k,m}^{\,t}(r) \in \gamma_i \text{ for some } t \ge 0 \,\}.
\]
By construction $\{0,\dots,m-1\} = \bigsqcup_{i=1}^{c} R_i$. The \textbf{modular weight} of $\gamma_i$ is $p_i = |R_i|/m = |R_i|/(b-1)$; note $\sum_i p_i = 1$.
\end{definition}

\begin{definition}[Residue signature]
For a physical attractor $A = \{a_1,\dots,a_L\} \in C(k,b)$, its \textbf{residue signature} is
\[
\sigma(A) = \{\, a_j \bmod m : 1 \le j \le L \,\} \subseteq \mathbb{Z}/m\mathbb{Z}.
\]
\end{definition}

\begin{definition}[Basin and density]
$B(A) = \{\, n \in \mathbb{Z}^+ : f_{k,b}^{\,t}(n) \in A \text{ for some } t \ge 0 \,\}$; its \textbf{natural density}, \emph{when the limit exists}, is $\delta(A) = \lim_{M\to\infty} |B(A) \cap [1,M]|/M$.
\end{definition}

\begin{definition}[Bifurcation excess]
$\Delta(k,b) = |C(k,b)| - \Cyc(\varphi_{k,b-1})$.
\end{definition}

\section{The Modular Skeleton}

\begin{lemma}[Modular invariance]\label{lem:inv}
For every $n \in \mathbb{Z}^+$, $f_{k,b}(n) \equiv n^k \pmod{m}$. Hence the residue sequence $\big(f_{k,b}^{\,t}(n) \bmod m\big)_{t\ge 0}$ is exactly the $\varphi_{k,m}$-orbit of $n \bmod m$.
\end{lemma}

\begin{proof}
By \eqref{eq:co9}, $f_{k,b}(n) = S_b(n^k) \equiv n^k \pmod m$, and induction on $t$ gives the orbit statement.
\end{proof}

\begin{lemma}[Contraction]\label{lem:contraction}
There is a threshold $N^*(k,b)$ such that $f_{k,b}(n) < n$ for all $n > N^*$. Consequently every orbit enters the finite set $[1, N^*]$ and is eventually periodic.
\end{lemma}

\begin{proof}
If $n$ has $D$ base-$b$ digits then $n \ge b^{D-1}$, while $n^k < b^{kD}$ has at most $kD$ digits, each at most $b-1$, so $S_b(n^k) \le (b-1)\,k\,D$. Since $b^{D-1}$ grows exponentially in $D$ and $(b-1)kD$ linearly, there is $D_0$ with $(b-1)kD < b^{D-1}$ for all $D \ge D_0$; take $N^* = b^{D_0 - 1}$.
\end{proof}

\begin{corollary}[Finite trapping region]\label{cor:trap}
Every attractor of $f_{k,b}$ lies in $[1, N^*(k,b)]$, and every basin $B(A)$ is determined by the first iterate landing in that window. (The miner of \S\ref{sec:verification} computes the least such fixed-point bound $M(k,b)$ and works entirely inside $[1,M]$.)
\end{corollary}

\section{The Lower Bound Theorem}\label{sec:lb}

\begin{theorem}[Lower Bound Theorem]\label{thm:lb}
For all $k \ge 1$ and $b \ge 2$,
\[
|C(k,b)| \;\ge\; \Cyc\big(\varphi_{k,\,b-1}\big).
\]
\end{theorem}

\begin{proof}
Three steps.

\textbf{Step 1 (Modular invariance).} By Lemma~\ref{lem:inv}, the residue class of $f_{k,b}(n)$ modulo $m$ is determined by that of $n$: the map $f_{k,b}$ respects the partition of $\mathbb{Z}^+$ into residue classes mod $m$, and the induced map on residues is exactly $\varphi_{k,m}$.

\textbf{Step 2 (Disjoint invariant subsets).} The cycles of $G(\varphi_{k,m})$ induce the basin partition $\{0,\dots,m-1\} = \bigsqcup_{i=1}^{c} R_i$. Set
\[
\mathcal{N}_i = \{\, n \in \mathbb{Z}^+ : n \bmod m \in R_i \,\}.
\]
Each $R_i$ is forward-invariant under $\varphi_{k,m}$ (if $r$ flows to $\gamma_i$, so does $\varphi(r)$), hence by Step~1 each $\mathcal{N}_i$ is forward-invariant under $f_{k,b}$, and the $\mathcal{N}_i$ are pairwise disjoint.

\textbf{Step 3 (One attractor per subset).} Fix $i$ and any $n_0 \in \mathcal{N}_i$. By Lemma~\ref{lem:contraction} the orbit of $n_0$ enters the finite set $[1,N^*] \cap \mathcal{N}_i$, which is forward-invariant by Step~2; being finite, the orbit eventually becomes periodic. The resulting periodic orbit is an attractor contained in $\mathcal{N}_i$.

\textbf{Conclusion.} The $\mathcal{N}_1,\dots,\mathcal{N}_c$ are pairwise disjoint and each contains at least one attractor, so $|C(k,b)| \ge c = \Cyc(\varphi_{k,b-1})$.
\end{proof}

\begin{remark}[Abstract form]\label{rem:abstract}
The proof uses only two properties of $f_{k,b}$: (i) a mod-$m$ semiconjugacy $f(n)\equiv\varphi(n\bmod m)$, and (ii) eventual contraction into a finite set. Thus for any eventually contracting $f:\mathbb{Z}^+\to\mathbb{Z}^+$ admitting a mod-$m$ semiconjugacy to a map $\varphi$ on $\mathbb{Z}/m\mathbb{Z}$, $|C(f)|\ge\Cyc(\varphi)$. In particular the bound applies to $f(n)=S_b(Q(n))$ for any $Q\in\mathbb{Z}[x]$. The happy maps $S_{e,b}$ fail (i): $\sum_i d_i^{\,e}$ is not a function of $n\bmod m$ alone (\S1.1).
\end{remark}

\begin{remark}
The bound is typically strict: in the verification of \S\ref{sec:verification} exact equality $|C| = \Cyc$ is the exception, not the rule. The excess $\Delta(k,b) \ge 0$ is a digit-dynamics effect (several attractors at different numerical scales sharing one residue signature) and is studied empirically in \S\ref{sec:split} and \S\ref{sec:predictors}.
\end{remark}

\begin{example}[The classic case $(k,b) = (2,10)$]\label{ex:classic}
The modular map $x \mapsto x^2 \bmod 9$ on $\{0,\dots,8\}$:
\[
0\to 0,\quad 1\to 1,\quad 2\to 4,\quad 3\to 0,\quad 4\to 7,\quad 5\to 7,\quad 6\to 0,\quad 7\to 4,\quad 8\to 1,
\]
with cycles $\{0\},\{1\},\{4,7\}$ and basins $R = \{0,3,6\},\{1,8\},\{2,4,5,7\}$. The physical attractors are $\{1\},\{9\},\{13,16\}$, exactly $3 = \Cyc(\varphi_{2,9})$, so $\Delta(2,10) = 0$: this is one of the rare exact-match systems.
\end{example}

\begin{example}[A bifurcating case $(k,b) = (3,10)$]\label{ex:bif}
The modular map on $\mathbb{Z}/9\mathbb{Z}$ has three fixed points $\{0\},\{1\},\{8\}$, predicting $\ge 3$ attractors. The physical system has seven: $\{1\},\{8\},\{17\},\{18\},\{19,28\},\{26\},\{27\}$. Within signature $\{0\}$ alone:
\[
9 \to S_{10}(729) = 18 \to S_{10}(5832) = 18 \ (\text{fixed}), \qquad 27 \to S_{10}(19683) = 27 \ (\text{fixed}).
\]
Both $18$ and $27$ are $\equiv 0 \pmod 9$ but are distinct fixed points at different scales: the modular theory counts one cycle for class $0$ while the real dynamics bifurcates. Hence $\Delta(3,10) = 7 - 3 = 4$.
\end{example}

\begin{proposition}[The case $k=1$]\label{prop:k1}
On $\mathbb{Z}^+$, for every $b\ge 2$, $C(1,b)=\{\{1\},\dots,\{b-1\}\}$, so $\Delta(1,b)=0$. Each individual basin $B(\{j\})$ equals $\{\,n\in\mathbb{Z}^+: n\equiv j\pmod{b-1}\,\}$ (taking $j=b-1$ for residue $0$) and has density $1/(b-1)$.
\end{proposition}

\begin{proof}
Write $n=\sum_{i<D} d_i b^i$ with $D\ge 2$ and $d_{D-1}\ge 1$. Then $n-S_b(n)=\sum_i d_i(b^i-1)\ge d_{D-1}(b^{D-1}-1)\ge b-1\ge 1$, so $S_b(n)<n$ for all $n\ge b$. Every cycle therefore lies in $[1,b-1]$, where $S_b$ is the identity. The modular map $\varphi_{1,b-1}$ is the identity on $\mathbb{Z}/(b-1)\mathbb{Z}$, hence $\Cyc=b-1=|C|$. The digital-root partition of $\mathbb{Z}^+$ is exactly the residue classes mod $b-1$, each of density $1/(b-1)$. (On $\mathbb{Z}_{\ge 0}$ the extra fixed point $0$ would give $\Delta=1$.)
\end{proof}

This is the clean extremal case of Remark~\ref{rem:scope}: one attractor per signature, no split, and the individual densities exist.

\section{The Basin Density Law (Aggregate Form)}\label{sec:bdl}

\begin{lemma}[Signature determinacy]\label{lem:sig}
For every $n \in \mathbb{Z}^+$, the attractor $A(n)$ eventually entered by the orbit of $n$ has residue signature
\[
\sigma\big(A(n)\big) = \gamma_i, \quad \text{where } i \text{ is the unique index with } (n \bmod m) \in R_i.
\]
\end{lemma}

\begin{proof}
By Lemma~\ref{lem:inv} the residue sequence of the orbit is the $\varphi_{k,m}$-orbit of $r = n \bmod m$, which enters the cycle $\gamma_i$ with $r \in R_i$. Once the physical orbit is on the attractor $A(n)$ (a cycle of period $L$), its residues are $\varphi$-periodic; travelling once around $A(n)$ returns to the start, so the residue period divides $L$, and since $\varphi$ restricted to a cycle is a cyclic permutation of that cycle, the residues traced out are the entire cycle $\gamma_i$. Hence $\sigma(A(n)) = \gamma_i$.
\end{proof}

\begin{proposition}[Exact finite-window mass partition]\label{prop:partition}
Fix $(k,b)$. For every $M \ge 1$,
\[
\sum_{\substack{A \in C(k,b) \\ \sigma(A) = \gamma_i}} \big|\, B(A) \cap [1,M] \,\big| \;=\; \big|\, \mathcal{N}_i \cap [1,M] \,\big|.
\]
\end{proposition}

\begin{proof}
The basins $\{B(A)\}_{A \in C(k,b)}$ partition $\mathbb{Z}^+$ (every orbit reaches a unique attractor, by Lemma~\ref{lem:contraction}). By Lemma~\ref{lem:sig}, $n \in \mathcal{N}_i$ iff $\sigma(A(n)) = \gamma_i$; thus the union of basins with signature $\gamma_i$ is exactly $\mathcal{N}_i$. Intersect with $[1,M]$ and count.
\end{proof}

\begin{theorem}[Basin Density Law, aggregate form]\label{thm:bdl}
For every $k \ge 1$, $b \ge 2$, and every modular cycle $\gamma_i$ of $\varphi_{k,b-1}$, the aggregate basin $\mathcal{N}_i = \bigcup_{\sigma(A) = \gamma_i} B(A)$ has a natural density, equal to the modular weight exactly:
\[
q_i \;:=\; \delta(\mathcal{N}_i) \;=\; p_i \;=\; \frac{|R_i|}{b-1}.
\]
Moreover, writing $M = qm + s$ with $0 \le s < m$, the empirical proportion $\hat q_i(M) = |\mathcal{N}_i \cap [1,M]|/M$ satisfies
\[
\big|\, \hat q_i(M) - p_i \,\big| \;\le\; \frac{\min\big(|R_i|,\; m-|R_i|,\; s\big)}{M} \;\le\; \frac{\min(|R_i|,\, b-2)}{M} \;\le\; \frac{b-1}{M}.
\]
\end{theorem}

\begin{proof}
By Proposition~\ref{prop:partition}, $\hat q_i(M) = |\mathcal{N}_i \cap [1,M]|/M$, and $\mathcal{N}_i \cap [1,M]$ is the set of integers in $[1,M]$ whose residue mod $m$ lies in $R_i$. Among $\{1,\dots,M\}$, exactly $s$ residue classes (call that set $S$) contain $q+1$ elements and the remaining $m-s$ contain $q$. Hence
\[
\big| \mathcal{N}_i \cap [1,M] \big| = |R_i| \cdot \frac{M}{m} + \theta, \qquad \theta = \sigma_i - \frac{s\,|R_i|}{m}, \quad \sigma_i = |R_i \cap S|.
\]
Since $\max\big(0, |R_i|-(m-s)\big) \le \sigma_i \le \min(|R_i|, s)$, evaluating $\theta$ at both endpoints gives $|\theta| \le \min(|R_i|, s)$. The deviations of all $m$ classes sum to zero, so the complement $\mathbb{Z}/m\mathbb{Z} \setminus R_i$ has deviation $-\theta$; applying the same estimate to it yields $|\theta| \le \min(m-|R_i|,\, s)$. Combining and dividing by $M$ gives the bound, and letting $M \to \infty$ yields $\delta(\mathcal{N}_i) = p_i$.
\end{proof}

\begin{remark}[The bound is not $(b-1)/M$]\label{rem:sharp}
The crude estimate $|\theta| \le |R_i| \le m$ gives $(b-1)/M$, which is what an earlier draft reported and described as saturated. It is not: since $s \le m-1$ and $\min(|R_i|, m-|R_i|) \le \lfloor m/2 \rfloor$, the constant $b-1$ is never attained, and over the sweep of \S\ref{sec:verification} the worst observed ratio of actual error to $(b-1)/M$ is $0.25$. The sharp constant $\min(|R_i|, b-2)$ is approached, with observed ratio $0.95$. The practical consequence is that a verification reporting ``$100\%$ of measurements within $(b-1)/M$'' tests very little; the statement with content is the integer identity of Proposition~\ref{prop:partition}, which has no slack.
\end{remark}

\begin{remark}[Scope of Theorem~\ref{thm:bdl}]\label{rem:scope}
The theorem asserts the density of the union $\mathcal{N}_i$. It does not assert that the individual basins $B(A)$ of the several attractors sharing signature $\gamma_i$ each have a natural density. When the local bifurcation excess is zero (one attractor per signature) the two coincide trivially. When it is positive, the decomposition $q_i = \sum_{\sigma(A)=\gamma_i} \delta(A)$ presupposes that the individual limits exist, which is false in general, as \S\ref{sec:split} shows: the individual densities oscillate and do not converge, while their sum remains exactly $p_i$. (A finite family of oscillating densities can sum to a set with exact density.) In the adjacent happy-number family the natural density of the happy set is likewise known not to converge (Gilmer \cite{gilmer}).
\end{remark}

\begin{corollary}[Conservation of mass between spaces]\label{cor:conservation}
Modular arithmetic determines the total basin mass of each residue signature exactly, independently of the bifurcation excess $\Delta(k,b)$. The excess governs only the internal partition of each mass $p_i$ among the $\ge 1$ physical attractors of signature $\gamma_i$, including whether that partition is well-defined as a vector of densities.
\end{corollary}

\begin{example}[The $22/33/44$ law]\label{ex:223344}
For $(k,b) = (2,10)$, Theorem~\ref{thm:bdl} with the basins of Example~\ref{ex:classic} predicts aggregate densities $3/9, 2/9, 4/9$ for signatures $\{0\},\{1\},\{4,7\}$, realized by the attractors $\{9\},\{1\},\{13,16\}$. Measured over $[1,3000]$: $33.33\% / 22.23\% / 44.43\%$, within the bound $9/3000 = 0.003$ of the exact values, as the theorem requires.
\end{example}

\section{Periodic Points and Cycles of the Modular Power Map}\label{sec:count}

Both counting functions of $\varphi_{k,m}(x) = x^k$ on $\mathbb{Z}/m\mathbb{Z}$ admit exact closed forms: the number of periodic points, which is classical, and the number of cycles, which is what Theorem~\ref{thm:lb} requires. Keeping them apart is essential; conflating the two produced an error in an earlier draft (Appendix~B.1).

\begin{proposition}[Periodic-point count]\label{prop:per}
For $k \ge 2$ and $m \ge 2$, the number of periodic points of $x \mapsto x^k$ on $\mathbb{Z}/m\mathbb{Z}$ is
\[
\#\Per(k,m) \;=\; \prod_{p^{e} \,\|\, m} \Big( 1 + \kappa_k\big(\varphi(p^{e})\big) \Big),
\]
where $\varphi$ is Euler's totient and $\kappa_k(N)$ is the largest divisor of $N$ coprime to $k$.
\end{proposition}

\begin{proof}
By the Chinese remainder theorem it suffices to treat $\mathbb{Z}/p^e\mathbb{Z}$. A nonzero non-unit is nilpotent under $x \mapsto x^k$ for $k \ge 2$: its $p$-adic valuation multiplies by $k$ at each step and reaches $\ge e$, so it flows to $0$; thus $0$ is the only periodic non-unit. On the unit group $(\mathbb{Z}/p^e\mathbb{Z})^\times$, of order $\varphi(p^e)$, an element $x$ is periodic under $g \mapsto g^k$ iff $x^{k^t} = x$ for some $t \ge 1$, iff $\ord(x) \mid k^t - 1$, iff $\gcd(\ord(x), k) = 1$. The set of such elements is the unique Hall subgroup of $k'$-torsion, whose order is the largest divisor of $\varphi(p^e)$ coprime to $k$, namely $\kappa_k(\varphi(p^e))$. Adding the point $0$ gives the local factor $1 + \kappa_k(\varphi(p^e))$; the factors multiply over primes by CRT.
\end{proof}

\begin{example}
For $(k,m) = (2,9)$: $\varphi(9) = 6$, $\kappa_2(6) = 3$, so $\#\Per = 1 + 3 = 4$, namely $\{0,1,4,7\}$, matching Example~\ref{ex:classic}.
\end{example}

The formula was checked against brute-force enumeration with 0 discrepancies across all $k \in [2,60]$, $m \in [2,259]$ (\texttt{scripts/cycle\_structure.py}), and independently re-verified in Appendix~A. It is classical in the theory of power digraphs $x \mapsto x^k \bmod n$ (Somer--K\v{r}\'i\v{z}ek \cite{somer-krizek}; Chou--Shparlinski \cite{chou-shparlinski}); it is included as the base for Proposition~\ref{prop:cyc}, not as a new result.

\begin{corollary}[Radical dependence of the periodic-point count]\label{cor:radical}
$\#\Per(k,m)$ depends on $k$ only through $\rad(k)$, the set of primes dividing $k$; in particular its $2$-part collapses entirely at the single transition $2 \nmid k \to 2 \mid k$ and is unchanged for all higher $2$-adic valuations $v_2(k)$.
\end{corollary}

\begin{proof}
Immediate from Proposition~\ref{prop:per}, since $\kappa_k(N)$ strips from $N$ exactly the primes dividing $k$.
\end{proof}

This property does not transfer to the cycle count. Corollary~\ref{cor:radical} concerns periodic points; Theorem~\ref{thm:lb} needs $\Cyc$. The two are not interchangeable, because a set of periodic points of known size can be organized into very different numbers of cycles. The correct statement on the cycle side is the following.

\begin{proposition}[Cycle count]\label{prop:cyc}
Let $k \ge 2$ and $m \ge 2$. For each prime power $p^e \,\|\, m$ let $H_{p^e} \le (\mathbb{Z}/p^e\mathbb{Z})^\times$ be the subgroup of elements whose order is coprime to $k$, and let $c_d$ be the number of elements of $H_{p^e}$ of order exactly $d$. The local cycle type of $\varphi_{k,p^e}$ on its periodic set is the multiset
\[
T_{p^e} \;=\; \{1\} \;\cup\; \bigcup_{d \,\mid\, \exp H_{p^e}} \Big\{\, \underbrace{\ord_d(k), \dots, \ord_d(k)}_{c_d/\ord_d(k) \ \text{times}} \,\Big\},
\]
where $\ord_d(k)$ is the multiplicative order of $k$ modulo $d$ (and $\ord_1(k) = 1$), and the singleton $\{1\}$ is the fixed point $0$. Then
\[
\Cyc(\varphi_{k,m}) \;=\; \sum_{(\ell_1,\dots,\ell_t) \,\in\, T_{p_1^{e_1}} \times \cdots \times T_{p_t^{e_t}}} \frac{\ell_1 \cdots \ell_t}{\lcmop(\ell_1,\dots,\ell_t)}.
\]
\end{proposition}

\begin{proof}
By Proposition~\ref{prop:per} the periodic set of $\varphi_{k,p^e}$ is $\{0\} \sqcup H_{p^e}$, and $\varphi_{k,m}$ restricted to it is a bijection. An element $x \in H_{p^e}$ of order $d$ satisfies $\varphi^{\,t}(x) = x^{k^t} = x$ iff $k^t \equiv 1 \pmod d$, so its cycle length is exactly $\ord_d(k)$; the $c_d$ elements of order $d$ therefore fill $c_d/\ord_d(k)$ cycles. By the Chinese remainder theorem the periodic set of $\varphi_{k,m}$ is the product of the local ones, and $\varphi$ acts coordinatewise. A tuple of local cycles of lengths $\ell_1,\dots,\ell_t$ spans $\prod \ell_i$ periodic points, on which the product permutation has order $\lcmop(\ell_i)$ and all orbits of that common length; hence it decomposes into $\prod \ell_i / \lcmop(\ell_i)$ cycles. Summing over tuples gives the formula.
\end{proof}

Checked against brute-force enumeration with 0 discrepancies for all $k \in [1,60]$, $m \in [1,259]$: $15{,}540$ pairs.

\begin{corollary}[The cycle count is not a function of $\rad(k)$]\label{cor:notradical}
$\Cyc(\varphi_{k,m})$ depends on $k$ through the multiplicative orders $\ord_d(k)$, hence on $k$ modulo the element orders $d$, not merely on the primes dividing $k$.
\end{corollary}

The dependence is not a marginal effect. For $m = 41$ the periodic set is $\{0\}$ together with the subgroup of order $5$, on which $x \mapsto x^k$ acts as multiplication by $k$ modulo $5$; since $\ord_5(2) = 4$, $\ord_5(4) = 2$ and $16 \equiv 1 \pmod 5$,
\[
\Cyc(\varphi_{2,41}) = 3, \qquad \Cyc(\varphi_{4,41}) = 4, \qquad \Cyc(\varphi_{16,41}) = 6,
\]
while $\#\Per = 6$ throughout. Likewise $\Cyc(\varphi_{2,37}) = 4$ but $\Cyc(\varphi_{4,37}) = 6$, with $\#\Per = 10$ in both cases. Across $m \in [2,500)$ there are $421$ moduli with $\Cyc(\varphi_{2,m}) \ne \Cyc(\varphi_{4,m})$, every one of them with $\#\Per$ identical, and $21$ of the $38$ moduli $m = b-1 \le 39$ inside the sweep grid of \S\ref{sec:verification} behave this way.

\begin{remark}[Scope]
Proposition~\ref{prop:cyc} closes the modular count side completely: both $\#\Per$ and $\Cyc$ are computable without building a graph. It says nothing about the physical excess $\Delta = |C| - \Cyc$, whose extra attractors are a digit-dynamics effect rather than a modular one; that remains open (Problem~\ref{prob:excess}). It also does not license any claim that $\Cyc$ is monotone or flat in $v_2(k)$; see Appendix~B.1 for the retracted statement.
\end{remark}

\section{Oscillation of the Intra-Signature Split}\label{sec:split}

Theorem~\ref{thm:bdl} fixes each aggregate mass $p_i$ exactly. The remaining question is how $p_i$ divides among the physical attractors sharing $\gamma_i$.

\subsection{Formulation}

Fix $(k,b)$ and a signature $\gamma_i$ hosting physical attractors $A_1,\dots,A_r$ ($r = 1 + {}$local excess), whose basins partition $\mathcal{N}_i$. Two questions:
\begin{itemize}
\item \textbf{(Existence)} Does each $\delta(A_j) = \lim_M |B(A_j) \cap [1,M]|/M$ exist?
\item \textbf{(Value)} If so, what is the vector $(\delta(A_1),\dots,\delta(A_r))$? By Proposition~\ref{prop:partition} it must satisfy $\sum_j \delta(A_j) = p_i$.
\end{itemize}

\subsection{Reduction to a digit-sum distribution}

By Corollary~\ref{cor:trap}, every attractor lies in $[1,M(k,b)]$, and one iterate already lands the orbit near that scale: $f_{k,b}(n) = S_b(n^k) \le (b-1)(1 + k \log_b n)$. Let $F = f_{k,b}\big|_{[1,M]}$ be the finite restricted map; each $A_j$ has a finite basin $\beta_j = B(A_j) \cap [1,M]$, and the $\beta_j$ partition $[1,M]$ within signature $\gamma_i$. Because basin membership of $n$ is decided by where its first iterate falls,
\[
B(A_j) = \{\, n : f_{k,b}^{\,t}(n) \in \beta_j \text{ for the first } t \text{ with } f^{\,t}(n) \le M \,\},
\]
so, tracking the first-passage landing site, the individual density would be
\[
\delta(A_j) \;=\; \sum_{v \in \beta_j} \dens\big\{\, n : \text{first iterate of } n \text{ into } [1,M] \text{ equals } v \,\big\},
\]
whenever the right-hand densities exist. The problem is thereby reduced to the asymptotic distribution of the digit sum $S_b(n^k)$ (and of its first few iterates) as $n$ ranges over a residue class, a well-studied object.

\subsection{Available analytic control}

The distribution of digit sums of polynomial sequences is a mature subject. For $k=2$, Mauduit and Rivat \cite{mauduit-rivat} proved $S_b(n^2)$ is asymptotically equidistributed in residue classes and satisfies a central limit theorem; Drmota--Mauduit--Rivat \cite{dmr} give local limit theorems for $S_b(n^k)$; Drmota \cite{drmota-poly} treats the summatory function. These control the pushforward measure above, and they also predict fluctuations, which turns out to be the essential point.

\subsection{Empirical resolution: the split oscillates}

Monte-Carlo measurement of the split restricted to $n$ with exactly $D$ base-$b$ digits (\texttt{scripts/split\_scale.py}, $(k,b) = (3,10)$, $4 \le D \le 90$, $120{,}000$ samples per band, sampling $\sigma \approx 0.0014$) gives a sharp, reproducible answer:
\begin{itemize}
\item The aggregate mass per signature is $1/3$ at every digit-length $D$ (Theorem~\ref{thm:bdl}, as expected).
\item No signature's split converges; all seven curves are non-monotone. For signature $\{0\}$ the two fixed points $\{18\}$ and $\{27\}$ exchange mass in antiphase, quasi-periodically in $D$, over the ranges $[0.060,0.239]$ and $[0.091,0.269]$, amplitude $\approx 0.179$, always summing to $1/3$. This is the largest and most visible instance.
\item The other two signatures oscillate as well, with smaller amplitude but the same qualitative behaviour. Within signature $\{1\}$ the $2$-cycle $\{19,28\}$ carries most of the mass but varies over $[0.266,0.338]$; within signature $\{8\}$ the fixed point $\{26\}$ varies over $[0.227,0.337]$.
\item The three small fixed points $\{1\},\{8\},\{17\}$ do not decay to zero. Each sits at or near zero for long stretches of $D$ and then returns: $\{17\}$ reaches $0.100$ at $D = 10$, is numerically zero throughout $22 \le D \le 70$, and comes back to $0.027$ at $D = 79$; $\{8\}$ is zero throughout $13 \le D \le 49$ and peaks at $0.028$ at $D = 58$; $\{1\}$ reaches $0.068$ at $D = 7$, vanishes over $13 \le D \le 64$, and returns to $0.026$ at $D = 73$. The returns are well above the sampling noise $\sigma \approx 0.0014$.
\end{itemize}

\begin{figure}[htbp]
\centering
\includesvg[width=0.95\linewidth]{../figures/split_oscillation}
\caption{Measured intra-signature split for $(k,b)=(3,10)$ as a function of digit length $D$. Competing attractors exchange mass in antiphase; the per-signature totals remain $1/3$.}
\label{fig:split-osc}
\end{figure}

\textbf{Mechanism.} In the reduction above the landing site is governed by the iterates $m_1 = S_b(n^k) \approx \tfrac{b-1}{2} k D$, which grows linearly in $D$, and then $m_2 = S_b(m_1^{\,k})$, whose mean grows only like $\log D$. Because the trapping region $[1,M]$ is a fixed window (here $M = 57$), the first-passage landing distribution lives on a fixed finite set and never escapes to infinity; it merely drifts and spreads slowly. As it sweeps across the fixed basin labelling of the integers, comparable basins alternate capture, and small basins sitting low in the window are revisited whenever the multi-step cascade produces a low landing value. The relevant basins inside $[1,57]$ are
\[
\beta_{\{1\}} = \{1,4,7,10,40\}, \quad \beta_{\{8\}} = \{2,5,8,11,20,50\}, \quad \beta_{\{17\}} = \{14,17,23,47\},
\]
each small but containing an element well above the others; those high elements are what the drifting landing distribution sweeps back across at large $D$. That is why the small fixed points recur instead of dying out.

\begin{remark}[Methodological note]
An earlier version of this subsection reported that signatures $\{1\}$ and $\{8\}$ converge. That conclusion came from a verdict rule which examined the drift of the split over a single step of $D$ and declared stabilization when it fell below the sampling noise. A quasi-periodic curve has long flat stretches, so such a rule reports convergence whenever the sampled window happens to lie inside one. The corrected measurement reports the whole curve and the amplitude over the full range of $D$, and never infers convergence from a local slope. See Appendix~B.2. A Gaussian-window extension of the \S\ref{sec:split} model to $D=300$ (no sampling of $n$) is reported in \S\ref{sec:split} and Conjecture~\ref{con:splitpilot}.
\end{remark}

\subsection{A parameter-free predictive model}

Two fit metrics appear below and should not be conflated: \textbf{short-band MAE} on $4 \le D \le 60$ (and on $8 \le D \le 64$ with $12{,}000$ samples per band) versus \textbf{bridge MAE} on the full measured range $D=4\ldots 90$ with $120{,}000$ samples per band (\texttt{scripts/bridge\_check.py}).

The oscillation is accounted for by a parameter-free model (\texttt{scripts/split\_predict.py}): model $m_1 = S_b(n^k)$ as Gaussian with mean $\tfrac{b-1}{2}L$ and variance $L\,\tfrac{b^2-1}{12}$ (the digit count $L$ being a computable mixture over the $D$-digit band; this is the local limit theorem of \cite{mauduit-rivat,dmr}), place the Gaussian on the \textbf{image lattice} $v \equiv r^k \pmod m$ for each residue $r$ feeding the signature (not $v \equiv r$; casting out nines forces $S_b(n^k) \equiv n^k \pmod m$ exactly), convolve with the exact integer-to-attractor labelling $a(v)$, and scale by $p_i$. For signature $\{0\}$ of $(3,10)$ this reproduces the measured split curve $\delta_j(D)$ to $\mathrm{MAE} = 0.003$ across $4 \le D \le 60$, with no fitted parameters. On the longer band $8 \le D \le 64$ with $12{,}000$ samples per band the same model gives $\mathrm{MAE} = 0.003$. Comparing the full measured curve to \texttt{predict\_split} on $D=4\ldots 90$ ($120{,}000$ samples per band; \texttt{scripts/bridge\_check.py}) gives mean per-attractor absolute error \textbf{$\approx 0.0017$}, at the measurement noise floor ($\approx 0.0014$). The oscillation is therefore the Gaussian sweep of the digit-sum $m_1$ across the fixed basin labelling of the integers, on the correct congruence classes of $m_1$.

On the feeding lattice for $(3,10)$, every sampled $v$ reaches the trapping region $[1,M]$ under at most \textbf{two} applications of $f_{3,10}$ (\texttt{first\_landing} in \texttt{src/dspm/predict.py}; scan to $v \le 10^6$ in \texttt{scripts/g4\_landing.py}). A proof for all $v$ is open.

An isolated Q-class laboratory (\texttt{src/dspm/qmaps.py}, records in \texttt{data/qclass/}) extends the pipeline to $f(n)=S_b(Q(n))$ with lattice $v \equiv Q(r) \pmod m$. The first non-monomial split, $Q(x)=1+3x+2x^2$ at $b=10$, shows the same antiphase oscillation (amplitude $\approx 0.27$ on $8 \le D \le 64$, Pearson $r=-1$ between the two competing $2$-cycles) with $\mathrm{MAE} = 0.006$ against the corrected model; an earlier $\mathrm{MAE} \approx 0.17$ on the same instance came from placing mass on $v \equiv r$ rather than $v \equiv Q(r)$.

The same Gaussian model, pushed to $4 \le D \le 300$ via \texttt{scripts/sweep\_label.py} (\texttt{data/split/label\_sweep\_k3\_b10\_sig0\_D300\_latest.md}), does not flatten: decade amplitude on $10 \le D \le 99$ is $0.192$ (full $\{\log_{10} D\}$ period), on $100 \le D \le 300$ is $0.132$, and cross-decade phase correlation is Pearson $r \approx 0.37$ ($n=201$ overlapping $\{\log_{10} D\}$ points between decades $10^1$ and $10^2$). The curves remain in exact antiphase (Pearson $-1$, algebraic). This is \textbf{model output}, not Monte Carlo measurement. Extending the diagnostic to $D=1000$ (\texttt{data/split/label\_sweep\_k3\_b10\_sig0\_D1000\_latest.md}) shows \textbf{amplitude survival} ($0.19$ on $10 \le D \le 99$; $0.13$ on $100 \le D \le 1000$), but cross-decade phase correlation on the overlapping $10^1$--$10^2$ interval weakens to $r \approx 0.26$. Long-scale phase organization is \textbf{not} claimed.

\begin{figure}[htbp]
\centering
\includesvg[width=0.95\linewidth]{../figures/split_predict_overlay}
\caption{Parameter-free Gaussian-sweep prediction (curves) against the Monte-Carlo split of Figure~\ref{fig:split-osc} for signature $\{0\}$, $(k,b)=(3,10)$. Short-band MAE $= 0.003$ on $4 \le D \le 60$; bridge MAE $\approx 0.0017$ on $D=4\ldots 90$ (image lattice $v \equiv r^k \pmod 9$).}
\label{fig:split-predict}
\end{figure}

With $120{,}000$ samples per band the measurement noise floor is $0.5/\sqrt{120000} \approx 0.0014$, so a bridge MAE of $0.0017$ says the model is consistent with the data at the noise floor. It does not single the model out against other candidates that also reproduce the phase and amplitude. The content of the model is mechanistic rather than statistical: it predicts the oscillation from the local limit theorem plus the exact labelling, with nothing tuned. Discriminating it from alternatives would need either a much larger sample or a prediction of the higher-order structure, which is what Problem~\ref{prob:split} asks for.

\begin{remark}[Empirical observation, mechanistically explained]\label{rem:result}
No individual basin density $\delta(A_j)$ appears to exist: along fixed digit-length the masses oscillate quasi-periodically in $D$, for every signature with a positive local excess. Only the per-signature aggregate has a density, equal to $p_i$ (Theorem~\ref{thm:bdl}).

This is recorded as an observation rather than a result: non-existence of a limit is a statement about all $D$, and finitely many bands cannot establish it. \textbf{Monte Carlo} (\texttt{split\_scale.py}, $4 \le D \le 90$, $120{,}000$ samples per band): the individual masses have not settled by $D=90$, with amplitude $\approx 0.18$ on signature $\{0\}$, well above the sampling noise $\sigma \approx 0.0014$. \textbf{Gaussian model} (\texttt{sweep\_label.py}, $4 \le D \le 300$, no sampling of $n$; \texttt{label\_sweep\_k3\_b10\_sig0\_D300\_latest.md}): persistent oscillation with decade amplitude $\approx 0.19$--$0.13$ and cross-decade phase correlation $\approx 0.37$, consistent with amplitude survival under the model (phase correlation is weaker than at small $D$). The step to the cumulative density is elementary: writing $\delta_j(d)$ for the mass among $d$-digit integers, the cumulative proportion up to $b^D$ is a weighted average with weights $\propto b^d$, so band $D$ alone carries weight $1-1/b$, and a persistent oscillation of amplitude $a$ survives with amplitude at least $(1-1/b)a$ minus the geometrically damped tail.
\end{remark}

\subsection{Relation to the theory of digital functions}

I have not found this basin-split oscillation written down verbatim for the $S_b(n^k)$ map. It is not a new kind of phenomenon: it is a manifestation of the classical log-periodic fluctuation of digital functions (Delange \cite{delange}, 1975), and the driving distribution $m_1 = S_b(n^k)$ is exactly the object of Drmota's work on the sum of digits of polynomial sequences \cite{drmota-poly} and of the Drmota--Grabner monograph \cite{drmota-grabner}. The $D=300$ Gaussian sweep shows amplitude survival with partial log-periodic organization (cross-decade phase correlation $\approx 0.37$ on $10^1$ vs $10^2$). The Gaussian-sweep model above is the standard Mellin--Perron / polynomial-digit-sum machinery of that theory, applied numerically to this map. A specialist would regard the oscillation as expected and the tools to make it rigorous as standard. What is claimed here is therefore: (i) an exact, elementary, computationally verified aggregate law (Theorem~\ref{thm:bdl}); (ii) an empirical and modelled description of why the individual densities fail to exist for this family, realized as an instance of digital-function oscillation. The next step is to check \cite{delange,drmota-poly,drmota-grabner} for whether the specific basin-split statement is already implied, and, if a note is warranted, to derive the exact Fourier/fractal oscillation term via Mellin--Perron rather than the numerical Gaussian.

\section{Computational Verification}\label{sec:verification}

\subsection{Method}

An exhaustive, finite-by-construction sweep (\texttt{scripts/sweep.py}). For each $(k,b)$ it computes the contraction bound $M(k,b)$, the fixed point of $S_b(n^k) \le (b-1)\cdot\mathrm{digits}_b(n^k)$, which by Lemma~\ref{lem:contraction} guarantees every attractor lies in $[1,M]$. It builds the full functional graph of $f_{k,b}$ on $[1,M]$, extracts all attractors, basins, and residue signatures $\sigma(A) \bmod (b-1)$; an independent engine reconstructs $G(\varphi_{k,b-1})$ and the weights $p_i$. Arbitrary precision via \texttt{gmpy2}; the sweep is parallelized over 28 processes.

\subsection{Results of the exhaustive sweep}

\begin{center}
\begin{tabular}{@{}ll@{}}
\toprule
\textbf{Parameter} & \textbf{Value} \\
\midrule
Exponents $k$ & $1$ to $500$ \\
Bases $b$ & $2$ to $40$ \\
Total pairs swept (exhaustive) & $19{,}500$ \\
Lower-bound violations ($|C| < \Cyc$) & $0$ \\
Failures of the exact integer identity (Prop.~\ref{prop:partition}) & $0$ \\
Failures of Lemma~\ref{lem:sig} (signature $=$ full modular cycle) & $0$ \\
Signature-density comparisons & $152{,}276$ \\
Points with error $\le \min(|R_i|,b-2)/M$ (sharp bound) & $152{,}276$ ($100.00\%$) \\
Mean absolute error $|\hat q_i - p_i|$ & $0.0001$ \\
Worst single error & $0.10$ at $(k,b,M) = (1,3,5)$ \\
Maximum observed excess & $\Delta = 98$ at $(k,b) = (451,32)$ \\
\bottomrule
\end{tabular}
\end{center}

The first three rows are the substantive ones: they are exact integer statements, so a single counterexample anywhere in the grid would falsify a theorem. The density rows are weaker than they look. The bound $(b-1)/M$ is not approached (over the grid the worst ratio of actual error to bound is $0.25$; Remark~\ref{rem:sharp}), so ``$100.00\%$ within the bound'' is close to automatic and should not be read as a stringent test. Reporting $152{,}276$ comparisons inflates the apparent evidence for the same reason: the comparisons are not independent, since all signatures of one pair share a window.

The sharp form is the integer identity of Proposition~\ref{prop:partition}: for every $(k,b)$ and every window $[1,M]$, the number of $n \le M$ whose orbit reaches an attractor of signature $\gamma_i$ equals, exactly, the number of $n \le M$ with $n \bmod (b-1) \in R_i$. No floating point is involved and there is no slack to absorb an error. That identity held without exception across the grid, and it is what the verification should be judged on.

This verifies the aggregate law; it says nothing about the split, which \S\ref{sec:split} resolves separately.

\subsection{Earlier sampled verification (superseded)}

An earlier prototype sweep tested $1{,}990$ sampled pairs (an incremental sample, not an exhaustive grid) with $k$ reaching $158$ and $b$ reaching $20$, using a finite sample $n \le 500$ per cell: zero violations, with exact equality $|C| = \Cyc$ in only $2$ of $1{,}990$ pairs ($\approx 0.1\%$). An independent denser re-sample over $k \in [1,25]$, $b \in [2,12]$ (275 pairs, $N = 300$) found 0 violations and a $7.64\%$ exact-match rate. The exact-match frequency is sensitive to the sampled range, itself an open question (Problem~\ref{prob:tight}). These figures are superseded by the exhaustive sweep above but are reported for provenance.

\subsection{Bifurcation heatmap phenomenology}

A systematic rendering of $|C(k,b)|$ over the $(k,b)$ grid shows:
\begin{enumerate}
\item \textbf{Vertical gradient:} larger bases produce more attractors, consistent with $\Cyc(\varphi_{k,b-1})$ growing with $b-1$.
\item \textbf{Columnar periodicity:} certain $k$ produce systematically fewer or more attractors across all bases, reflecting the structure of $\gcd(k, b-1)$.
\item \textbf{Diagonal resonances:} hot spots follow diagonal trajectories, suggesting dependence on $k \bmod (b-1)$ or $\gcd(k, \varphi(b-1))$.
\end{enumerate}

\section{Structural Predictors of the Attractor Count (Empirical)}\label{sec:predictors}

Theorem~\ref{thm:bdl} fixes the $p_i$ but is silent on how many physical attractors share $\gamma_i$. I report empirical predictors of the count side (the excess $\Delta$), stated as correlations, not closed forms, with flags where they overlap known happy-number structure.

\medskip
\noindent\textbf{9.1 $\omega(b-1)$ is the master structural predictor.} With $\omega(m)$ the number of distinct prime factors of $m = b-1$: Pearson $R = +0.436$ with $\Cyc$, $+0.419$ with $|C|$, $+0.391$ with $\Delta$. To maximize macro-states, take $b-1$ with many distinct primes ($b-1 = 2\cdot3\cdot5 = 30 \Rightarrow b = 31$).

\medskip
\noindent\textbf{9.2 The $2$-adic structure of $\gcd(k, b-1)$ controls the excess in a single step.} The linear correlation of $\gcd$ with $\Delta$ is negligible ($-0.05$), masking a strong pattern under grouping:

\begin{center}
\begin{tabular}{@{}cc@{\qquad}cc@{}}
\toprule
$\gcd(k,b-1)$ & mean $\Delta$ & $\gcd$ & mean $\Delta$ \\
\midrule
$1$ & $14.46$ & $3$ & $14.15$ \\
$2$ & $8.79$ & $5$ & $16.55$ \\
$4$ & $6.90$ & $15$ & $26.88$ (max) \\
$8$ & $5.16$ & & \\
$16$ & $4.06$ & & \\
$32$ & $3.73$ (min) & & \\
\bottomrule
\end{tabular}
\end{center}

Powers of two lower $\Delta$; odd gcds (especially with several odd primes) inflate it.

\textbf{Caution on reading this table.} Each row pools pairs from many different moduli, and $v_2(\gcd(k,b-1))$ is large only when $b-1$ itself carries $2$-power structure, so the rows are not comparable populations. Within a single modulus the cycle count is neither monotone nor flat in $v_2(k)$: by Corollary~\ref{cor:notradical}, $m=41$ gives $\Cyc = 3, 4, 3, 6, 3$ for $k = 2, 4, 8, 16, 32$. The decreasing trend down the powers-of-two column is therefore an aggregation effect across moduli, and no per-system law should be read off it. \texttt{scripts/analyze\_patterns.py} reports the within-group standard deviation and the number of distinct moduli in each row for exactly this reason. An earlier draft asserted a graded geometric collapse $\overline{\Delta}(v_2) \approx 12.68 \cdot (0.764)^{v_2}$; a later one replaced it with the claim that the cycle count is flat for all $v_2 \ge 1$. Both are withdrawn; see Appendix~B.1.

\textbf{Relation to known results.} This is not independent of the happy-number literature: Grundman--Teeple \cite{GT07} show arbitrarily long runs of consecutive $(e,b)$-happy numbers exist precisely when $e-1$ is not divisible by $p-1$ for any prime $p \mid b-1$, a coprimality condition on the exponent against the prime structure of $b-1$ of the same flavour as the pattern above. Sections 9.1--9.2 are an empirical, quantitative counterpart for the digit-sum-power family, not a claim independent of that circle of ideas.

\medskip
\noindent\textbf{9.3 The parity of $k$, not its primality, amplifies attractors.} The raw contrast replicates (prime $k$: mean $|C| = 31.82$ against $18.26$ for composite, about $74\%$ more) but it is confounded. Two is the only even prime, while roughly half of the composites are even, and $2 \mid k$ collapses the $2$-part of every local unit group (Proposition~\ref{prop:per}: $\kappa_k$ strips it), cutting both $\Cyc$ and $|C|$. Controlling for parity on an exhaustive grid of $1{,}298$ pairs ($2 \le k \le 60$, $3 \le b \le 24$):

\begin{center}
\begin{tabular}{@{}lccc@{}}
\toprule
class of $k$ & $n$ & mean $|C|$ & mean $\Cyc$ \\
\midrule
prime (raw) & $374$ & $17.27$ & $7.97$ \\
composite (raw) & $924$ & $11.40$ & $4.82$ \\
odd prime & $352$ & $18.09$ & $8.25$ \\
odd composite & $286$ & $17.16$ & $7.35$ \\
odd (any) & $638$ & $17.67$ & $7.85$ \\
even (any) & $660$ & $8.66$ & $3.68$ \\
\bottomrule
\end{tabular}
\end{center}

Odd primes and odd composites differ by $0.93$ in mean $|C|$, against a within-class standard deviation of about $9.7$; odd and even differ by $9.01$. The effect is parity. The related correlation $R = -0.391$ between the number-of-divisors function of $k$ and $\Cyc$ is subject to the same confound, since even $k$ tends to have more divisors. (Appendix~B.4.)

\medskip
\noindent\textbf{9.4 Transient depth scales with $k$.} Maximum transient depth correlates $R = +0.540$ with $k$; deepest observed: $133$ steps.

\section{Open Problems}\label{sec:open}

\begin{problem}[Exact formula]
Find a closed-form expression or efficient algorithm for $|C(k,b)|$ beyond Theorem~\ref{thm:lb}. In particular, characterize the bifurcation excess $\Delta(k,b)$.
\end{problem}

\begin{problem}[Tightness characterization]\label{prob:tight}
For which pairs $(k,b)$ does equality $|C(k,b)| = \Cyc(\varphi_{k,b-1})$ hold? For $k=1$ equality holds for every $b\ge 2$ (Proposition~\ref{prop:k1}). Data otherwise suggests small $b$ (especially $b = 2,3$) or trivial modular dynamics; the exact-match frequency is itself range-sensitive (\S\ref{sec:verification}). A complete census for $b=2,3$ and $1 \le k \le 1500$ finds equality precisely at $k \in \{1,2,3,7,381\}$ when $b=2$ and at $k \in \{1,6,10\}$ when $b=3$; no further exact matches appear after $k=400$. Outside $k=1$ this is a list, not a characterization.
\end{problem}

\begin{problem}[Orbit length asymptotics]
Let $L(k,b,N)$ be the maximum orbit length (transient + cycle) among $n \le N$. What is its growth rate? Is $L(k,b,N) = O(\log N)$ for all fixed $(k,b)$?
\end{problem}

\begin{problem}[Predecessor graph structure]
For fixed $(k,b)$ and $A \in C(k,b)$, the basin tree rooted at $A$: what is its degree distribution? Does it follow a power law? On the functional graph of $f_{k,b}$ restricted to $[1,M]$, a Clauset--Shalizi--Newman discrete power-law fit is not plausible for any of the systems checked.
\end{problem}

\begin{problem}[Upper bound]
Prove a complementary upper bound on $|C(k,b)|$; a natural candidate is the contraction threshold $N^*(k,b)$ of Lemma~\ref{lem:contraction}. The inequality $|C| \le N^*$ holds and is uninformative: $N^*$ bounds attractor \emph{values}, not their count. Empirically, $\Cyc$ times the number of attractor digit layers already fails as an upper bound on most pairs of a stratified follow-up grid; no sharp complementary inequality survived.
\end{problem}

\begin{problem}[Rigorous theory of the split]\label{prob:split}
Prove that the per-digit split does not converge when a signature splits; determine whether non-convergence is already implied by \cite{delange,drmota-poly,drmota-grabner}. No single period-1 Delange factor $P_j(\{\log_b D\})$ is claimed. The target statements are as follows. Nothing below is claimed as a theorem.

\medskip\noindent\textbf{Setup.} Fix $k\ge 2$, $b\ge 2$, and write $m=b-1$, $M=M(k,b)$ for the contraction bound of Lemma~\ref{lem:contraction}. Let $N_D=\{n: b^{D-1}\le n<b^D\}$. For $v\ge 1$ let $g(v)$ be the first iterate of $f_{k,b}$ that lies in $[1,M]$ (so $g(v)=v$ if $v\le M$). For an attractor $A_j$ write $\beta_j=B(A_j)\cap[1,M]$ and
\[
\delta_j(D)\;=\;\frac1{|N_D|}\,\bigl|\{\,n\in N_D: g\bigl(S_b(n^k)\bigr)\in\beta_j\,\}\bigr|.
\]
This is the per-digit split of \S\ref{sec:split}. The map $g$ is completely explicit: it is determined by the finite functional graph on $[1,M]$. Theorem~\ref{thm:bdl} already gives $\sum_{A_j\subset\gamma_i}\delta_j(D)\to p_i$ as $D\to\infty$ (in fact the finite-window identity of Proposition~\ref{prop:partition} controls every prefix). The remaining claim is that no individual $\delta_j(D)$ converges when $\gamma_i$ splits.

\medskip\noindent\textbf{Hypothesis LLT$(k,b,r)$.} For each residue $r\bmod m$, writing $N_D^{(r)}=N_D\cap(r+m\mathbb{Z})$ and taking $n$ uniform in $N_D^{(r)}$, the law of $S_b(n^k)$ admits a local Gaussian limit with mean $\mu_D\asymp D$ and standard deviation $\sigma_D\asymp\sqrt{D}$: for every fixed $C<\infty$,
\[
\sup_{\lvert v-\mu_D\rvert\le C\sigma_D}\Bigl\lvert \sigma_D\,\mathbb{P}\bigl(S_b(n^k)=v\bigr)-\varphi\bigl((v-\mu_D)/\sigma_D\bigr)\Bigr\rvert\;\longrightarrow\;0
\]
as $D\to\infty$, and the tails $\mathbb{P}(\lvert S_b(n^k)-\mu_D\rvert>C\sigma_D)$ vanish as $C\to\infty$ uniformly in $D$. (Casting out nines forces $S_b(n^k)\equiv n^k\pmod m$ exactly.) This is the form in which a restriction of the local limit theorem of Drmota--Mauduit--Rivat \cite{dmr} to a dyadic band and an arithmetic progression would be used. It is an input, not a claim of the present paper.

\medskip\noindent\textbf{Hypothesis LM$(k,b,j)$.} Let $h_j(v)=1_{g(v)\in\beta_j}$. Write $\Psi_j(V)$ for the mean of $h_j$ on the window $[V-\sqrt{V},\,V+\sqrt{V}]$ restricted to integers $v \equiv r^k \pmod m$ on the image lattice of residues $r$ feeding the signature of $A_j$. The hypothesis is that $\Psi_j(V)$ does not converge as $V\to\infty$. This is a statement about the labelling, parallel to Hypothesis LLT and independent of it. It is diagnosed computationally by \texttt{scripts/local\_mean.py}; it is not deduced from a Delange expansion of $h_j$.

\medskip\noindent\textbf{Remark (thin window; retracted).} An earlier draft of this section asked to prove that if $\sum_{n\le x}f(n)=x\,P(\log_b x)+o(x)$ with $P$ continuous of period $1$, and if $I_x\subset[c_1 x,\,c_2 x]$ satisfies $|I_x|/x\to 0$, then the window mean equals $P(\log_b x)+o(1)$. That implication is false for every remainder rate. If $P$ is $C^1$ of period $1$, $h=|I_x|\to\infty$, $h=o(x)$, and $\sup|E|=o(h)$ on the window, the mean is $Q(\log_b x)+o(1)$ with $Q=P+P'/\ln b$, not $P$; $Q$ is constant if and only if $P$ is. Counterexample: in base $10$, $f(v)=1_{\text{leading digit of }v=1}$ has remainder $O(1)$; along $x=1.5\cdot 10^t$ one has $P(\log_{10}x)\approx 0.4074$ while the $\sqrt{x}$-window average is $1$, matching $Q=1$. Delange's periodic function for $s_b$ is continuous and nowhere differentiable, so the $C^1$ hypotheses fail in the intended regime and no such $Q$ should be expected. The lemma is withdrawn; its role is taken by Hypothesis LM.
\end{problem}

\begin{conjecture}\label{con:split}
Assume Hypothesis LLT$(k,b,r)$ for every residue $r$ feeding a split signature $\gamma_i$ (that is, $a_i\ge 2$), and Hypothesis LM$(k,b,j)$ for each attractor $A_j$ of that signature. Then $\lim_{D\to\infty}\delta_j(D)$ does not exist. The same non-existence passes to the cumulative density $|B(A_j)\cap[1,b^D]|/b^D$, because band $D$ carries geometric weight $1-1/b$ (\S\ref{sec:split}). No identity $\delta_j(D)=P_j(\{\log_b D\})+o(1)$ is claimed: under the Gaussian model the decade-collapse $\delta_j(D)\approx\delta_j(bD)$ already fails for the pilot of Conjecture~\ref{con:splitpilot}, consistently with the absence of a $C^1$ Delange local mean.
\end{conjecture}

\begin{conjecture}[Pilot]\label{con:splitpilot}
Take $(k,b)=(3,10)$ and the signature $\gamma=\{0\}$, whose physical attractors are $\{18\}$ and $\{27\}$ (Example~\ref{ex:bif}). Then $\delta_{18}(D)+\delta_{27}(D)\to 1/3$, yet neither factor converges: $\limsup\delta_{18}-\liminf\delta_{18}>0$, and the two curves are in antiphase. This is the content of Remark~\ref{rem:result} for the largest split, upgraded from a finite-$D$ measurement to a statement about all $D$.

A Gaussian-window diagnostic (\texttt{scripts/sweep\_label.py}; \texttt{data/split/label\_sweep\_k3\_b10\_sig0\_D300\_latest.md}) evaluates the \S\ref{sec:split} model on signature $\{0\}$ for $4 \le D \le 300$ without sampling $n$ (ceiling $V=4789$). On attractor $\{18\}$, decade amplitude is $0.192$ on $10 \le D \le 99$ (full $\{\log_{10} D\}$ period) and $0.132$ on $100 \le D \le 300$. Cross-decade phase correlation is Pearson $r \approx 0.37$ ($n=201$ overlapping $\{\log_{10} D\}$ points between decades $10^1$ and $10^2$). The curves $\delta_{18}(D)$ and $\delta_{27}(D)$ remain in exact antiphase (Pearson $-1$, algebraic), always summing to $1/3$. Under Hypothesis LLT, \textbf{amplitude survival} in this range is consistent with Conjecture~\ref{con:splitpilot}; strong cross-decade phase locking is \textbf{not} claimed. Extending to $D=1000$ (\texttt{label\_sweep\_k3\_b10\_sig0\_D1000\_latest.md}) amplitude persists but phase correlation weakens to $r \approx 0.26$ (marked \textbf{amplitude\_only} in the diagnostic pipeline).
\end{conjecture}

The Gaussian of \S\ref{sec:split} is the leading LLT term. Comparing MAE of a truncated Fourier inversion against the Gaussian does not prove the conjecture (the two already agree to sampling noise). The work is: (i) Hypothesis LM for the labelling, diagnosed by the $v$-space means $\Psi_j(V)$ (\texttt{scripts/local\_mean.py}); (ii) the large-$D$ Gaussian diagnostic (\texttt{scripts/sweep\_label.py}); (iii) whether \cite{delange,drmota-poly,drmota-grabner} imply LM; (iv) the LLT input.

\begin{problem}[Count side of the excess]\label{prob:excess}
The modular count side is now fully closed: both the periodic-point count (Proposition~\ref{prop:per}) and the cycle count (Proposition~\ref{prop:cyc}) have closed forms. Characterize the physical excess $\Delta = |C| - \Cyc$, extra attractors from digit dynamics rather than modular structure. Empirically $\Delta$ responds to the parity of $k$ and to $\omega(b-1)$ (\S\ref{sec:predictors}), but no closed form is known.
\end{problem}

\begin{problem}[Complexity of the cycle count; resolved]
Proposition~\ref{prop:cyc} evaluates $\Cyc(\varphi_{k,m})$ from the factorization of $m$ and multiplicative orders modulo divisors of $\lambda(m)$, so it is polynomial given the factorization; the naive sum over $t$-tuples of local cycles is not. The same quantity is obtained by folding multiplicity maps of cycle lengths: a cycle of length $\ell_1$ and a cycle of length $\ell_2$ produce $\gcd(\ell_1,\ell_2)$ cycles of length $\lcmop(\ell_1,\ell_2)$. Pairwise CRT folding never enumerates $t$-tuples (\texttt{cycle\_count\_formula\_folded} in \texttt{dspm.modular}; equal to the graph and to the expanded product in \texttt{tests/test\_modular.py}). Cost is polynomial in the size of the local cycle-type maps, given the factorization.
\end{problem}

\section{Conclusion}

For the family $f_{k,b}(n) = S_b(n^k)$, modular arithmetic modulo $b-1$ determines, exactly and elementarily, both a lower bound on the number of attractors (Theorem~\ref{thm:lb}) and the total basin mass of each residue signature (Theorem~\ref{thm:bdl}, whose finite-window content is an exact identity between integers, with the sharp bound $\min(|R_i|,b-2)/M$). Both were verified without exception across the exhaustive sweep of $19{,}500$ systems. The modular count side is closed in both of its counting functions: the classical periodic-point formula of Proposition~\ref{prop:per}, and the cycle-count formula of Proposition~\ref{prop:cyc}, which is the one Theorem~\ref{thm:lb} needs, and which shows that $\Cyc$, unlike $\#\Per$, is not a function of $\rad(k)$.

These exact results quotient out the elementary part of the distribution and isolate the remaining object: the intra-signature split. That split does not settle. Along fixed digit-length the individual basin masses oscillate quasi-periodically, in antiphase between competing attractors, always summing to the exact aggregate $p_i$ (Remark~\ref{rem:result}). This holds for every signature with a positive local excess, including the ones an earlier draft classified as convergent. A parameter-free Gaussian-sweep model reproduces the oscillation to within Monte-Carlo noise. The oscillation is an instance of the classical log-periodic fluctuation of digital functions (Delange; Drmota--Grabner). The contribution of the present work is a correctly stated, computationally verified aggregate law, a closed form for the cycle count, and a characterization of why the finer distribution appears not to exist for this family.

\appendix

\section{Reproducibility and Independent Verification}

\textbf{Primary pipeline.} Everything is reproducible from the \texttt{dspm} package accompanying this article. The exhaustive sweep behind \S8.2 is \texttt{scripts/sweep.py}, producing \texttt{data/sweeps/results\_k1-500\_b2-40\_*.jsonl.gz} (19,500 records). The split programme of \S\ref{sec:split} is \texttt{scripts/split\_scale.py} (measurement) and \texttt{scripts/split\_predict.py} (the parameter-free model); bridge error is \texttt{scripts/bridge\_check.py}; the $v$-space local means of Hypothesis LM are \texttt{scripts/local\_mean.py}; the large-$D$ Gaussian diagnostic of Conjecture~\ref{con:splitpilot} is \texttt{scripts/sweep\_label.py}. LM sidecar diagnostics: \texttt{scripts/lm\_stratum.py}, \texttt{scripts/lm\_oscillation.py}, \texttt{scripts/lm\_suffix.py}, \texttt{scripts/lm\_carry\_depth.py}, \texttt{scripts/g4\_landing.py} (outputs in \texttt{data/qclass/split/}). An isolated Q-class laboratory for $S_b(Q(n))$ lives in \texttt{src/dspm/qmaps.py} with records in \texttt{data/qclass/}. The count side of \S\ref{sec:count} is \texttt{scripts/cycle\_structure.py}. Statistics of \S\ref{sec:predictors} are \texttt{scripts/analyze\_patterns.py}. The computational remarks in \S\ref{sec:open} (tightness census, predecessor degrees, bound candidates, folded cycle count) are \texttt{scripts/analyze\_topic10.py} and \texttt{dspm.mining}. Figures: \texttt{scripts/plot\_split\_figures.py} writes \texttt{paper/figures/split\_oscillation.svg} and \texttt{split\_predict\_overlay.svg}.

\begin{verbatim}
pip install -e ".[fast,dev]"
pytest                                        # theorem-level test suite
python verification/verify_theorems.py        # independent re-derivation
python scripts/cycle_structure.py             # both closed forms vs brute force
python scripts/split_scale.py --k 3 --b 10 --d-max 90 --samples 120000
python scripts/bridge_check.py --k 3 --b 10
python scripts/sweep_label.py --k 3 --b 10 --d-max 300
python scripts/sweep_label.py --k 3 --b 10 --d-max 1000
python scripts/local_mean.py --k 3 --b 10 --v-max 10000000
python scripts/plot_split_figures.py --k 3 --b 10 --d-max 90 --lang en --oscillation
\end{verbatim}

\medskip
\noindent\textbf{Independent re-verification.} \texttt{verification/verify\_theorems.py} reimplements the digit sum, the trapping region, both functional graphs and the modular structure from scratch, importing nothing from \texttt{dspm}, so a bug in the package cannot make a theorem appear to hold. On $b \in [2,11]$, $k \in [1,15]$ it reports:
\begin{itemize}
\item Lemma~\ref{lem:inv} (modular invariance): $40{,}365$ tests, 0 failures.
\item Theorem~\ref{thm:lb} ($|C| \ge \Cyc$): exhaustive over $150$ pairs, 0 violations; exact equality in $19$ of them.
\item Lemma~\ref{lem:sig} (signature is a full modular cycle): 0 violations.
\item Proposition~\ref{prop:partition} (exact integer identity), over three windows per pair: 0 violations.
\item Theorem~\ref{thm:bdl}, against both the stated bound and the sharp $\min(|R_i|,b-2)/M$: 0 violations; worst error-to-bound ratio $0.25$ for the former and $0.90$ for the latter, which is the evidence that the original bound was not sharp.
\item Proposition~\ref{prop:per}: closed form against brute force, $7{,}722$ pairs, 0 discrepancies.
\item $(k,b)=(2,10)$: attractors $\{1\},\{9\},\{13,16\}$ with measured basin proportions $0.2223/0.3333/0.4443$ on $[1,3000]$, matching $2/9, 3/9, 4/9$.
\item $(k,b)=(3,10)$: attractors $\{1\},\{8\},\{17\},\{18\},\{19,28\},\{26\},\{27\}$, matching Example~\ref{ex:bif}.
\end{itemize}

Proposition~\ref{prop:cyc} is checked separately against brute-force enumeration for all $k \in [1,60]$, $m \in [1,259]$: $15{,}540$ pairs, 0 discrepancies.

\section{Errata}\label{app:errata}

Four claims in the July 2026 draft were wrong. Each is recorded with the reasoning that produced it, because in three of the four cases the error was in the verification method rather than in the mathematics, and those failure modes generalize.

\medskip
\noindent\textbf{B.1. The modular cycle count is flat for $v_2(k) \ge 1$.} \emph{Withdrawn.} The claim was stated as a consequence of Proposition~\ref{prop:per} and ``verified for $m = 16, 17, 32, 37, 41, 64$''. Two independent defects. First, Proposition~\ref{prop:per} counts periodic points, and $\#\Per$ is indeed a function of $\rad(k)$; the cycle count is not, because cycle lengths are multiplicative orders $\ord_d(k)$ (Proposition~\ref{prop:cyc}, Corollary~\ref{cor:notradical}). Second, the supporting computation grouped $k$ by the odd part and $v_2$ of $\gcd(k,\lambda(m))$ and compared group means. For the two non-degenerate moduli in the list the within-group values are $\{4,6,10\}$ for $m=37$ and $\{3,4,6\}$ for $m=41$; the means agree across $v_2$ groups because each group contains the same multiset in the same proportions. The other four moduli are degenerate: $\kappa_k$ collapses the unit part, $\#\Per = 2$, and flatness is vacuous. The lesson is that a test comparing group means cannot detect non-constancy, and that four of six test cases were incapable of failing. Replaced by Proposition~\ref{prop:cyc} and Corollary~\ref{cor:notradical}.

\medskip
\noindent\textbf{B.2. Signatures $\{1\}$ and $\{8\}$ have convergent splits.} \emph{Withdrawn.} The measurement script decided convergence from the change in the split across one step of $D$, calling it stabilized when that change fell below the sampling noise. Because the curves are quasi-periodic with long flat stretches, the rule reports convergence whenever the window sampled lies inside one, and the original run stopped at $D = 60$, inside such a stretch for two of the three fixed points. Re-measuring on $4 \le D \le 90$ with $60{,}000$ samples per band shows all three ``vanishing'' fixed points returning after long silences, with peaks of $0.027$ at $D = 79$ for $\{17\}$, $0.028$ at $D = 58$ for $\{8\}$, and $0.026$ at $D = 73$ for $\{1\}$; all seven curves are non-monotone.\footnote{The current Tier~B measurement uses $120{,}000$ samples per band on $4 \le D \le 90$; the $60{,}000$ figure records the August 2026 erratum re-measurement.} Corrected in \S7.4. The correction strengthens the main empirical claim: every individual density fails to settle, not merely the one in signature $\{0\}$.

\medskip
\noindent\textbf{B.3. The finite-window bound $(b-1)/M$ is ``saturated''.} \emph{Withdrawn.} It is not: over the grid the worst ratio of actual error to bound is $0.25$. The proof of Theorem~\ref{thm:bdl} bounds the deviation of one residue class by $|R_i|$, but the deviations across all $m$ classes sum to zero, so the sharp constant is $\min(|R_i|,\,m-|R_i|)$, and in a window $M = qm + s$ it is $\min(|R_i|, s) \le \min(|R_i|, b-2)$. Section~8.2 already displayed the gap (worst error $0.10$ against a bound of $0.40$) without drawing the conclusion. Corrected in \S1.3, \S8.2 and \S11; the sharp bound is verified in \texttt{verification/audit/audit\_03\_bound\_sharpness.py}.

\medskip
\noindent\textbf{B.4. Prime exponents amplify attractors by $\approx 74\%$.} \emph{Reinterpreted.} The number replicates but the contrast is confounded with the parity of $k$, which is the actual driver (Proposition~\ref{prop:per}). Controlled comparison in \S9.3.

\medskip
\noindent\textbf{B.5. The Gaussian lattice used feeding residues of $n$.} \emph{Corrected.} An earlier implementation of \S\ref{sec:split} placed the Gaussian on $v \equiv r \pmod m$ for each residue $r$ feeding the signature. Casting out nines forces $S_b(n^k) \equiv n^k \pmod m$, so the correct lattice is $v \equiv r^k \pmod m$. On signature $\{0\}$ of $(3,10)$ the wrong lattice gave $\mathrm{MAE} \approx 0.041$ on the band $8 \le D \le 64$; the image lattice restores $\mathrm{MAE} \approx 0.003$. For a general polynomial $Q$, the sidecar uses $v \equiv Q(r) \pmod m$ (\texttt{predict\_split\_Q} in \texttt{src/dspm/qmaps.py}). Corrected in \texttt{predict\_split} and documented in \S\ref{sec:split}; diagnostic records in \texttt{data/qclass/split/twostep\_latest.md}.

\medskip
\noindent\textbf{A methodological note.} Four of the five errors survived because the verification was run against a statistic that could not distinguish the claim from its negation: a comparison of means where constancy was at issue, a local slope where a limit was at issue, a loose inequality where sharpness was at issue, and a congruence lattice where the pushforward of $m_1$ was at issue. The countermeasure adopted throughout the accompanying code is to prefer exact integer identities where they exist, as in Proposition~\ref{prop:partition}, and, where they do not, to report amplitude, spread and noise floor next to every verdict.

\begin{thebibliography}{99}

\bibitem{guy} R.~K. Guy, ``Happy Numbers,'' in \emph{Unsolved Problems in Number Theory}, 3rd ed., Springer, 2004, pp.~358--359.

\bibitem{GT01} H.~G. Grundman and E.~A. Teeple, ``Generalized happy numbers,'' \emph{The Fibonacci Quarterly} \textbf{39} (2001), 462--466.

\bibitem{GT07} H.~G. Grundman and E.~A. Teeple, ``Sequences of consecutive happy numbers,'' \emph{Rocky Mountain J. Math.}; and ``Sequences of generalized happy numbers with small bases,'' \emph{J. Integer Sequences} \textbf{10} (2007), Article 07.1.8.

\bibitem{hardy-wright} G.~H. Hardy and E.~M. Wright, \emph{An Introduction to the Theory of Numbers}, 6th ed., Oxford University Press, 2008.

\bibitem{survey} E.~Hasson et al.\ (eds.), ``Happy Numbers, Happy Functions, and Their Variations: A Survey,'' \emph{La Matematica} (Springer), 2021.

\bibitem{gilmer} J.~Gilmer, ``On the density of happy numbers,'' \emph{Integers} \textbf{13} (2013), \#A48 (arXiv:1110.3836).

\bibitem{mauduit-rivat} C.~Mauduit and J.~Rivat, ``La somme des chiffres des carr\'es,'' \emph{Acta Mathematica} \textbf{203} (2009), 107--148.

\bibitem{dmr} M.~Drmota, C.~Mauduit, and J.~Rivat, ``The sum of digits of polynomial values in arithmetic progressions,'' and related local-limit results for $S_b(n^k)$.

\bibitem{silverman} J.~H. Silverman, \emph{The Arithmetic of Dynamical Systems}, GTM 241, Springer, 2007.

\bibitem{somer-krizek} L.~Somer and M.~K\v{r}\'i\v{z}ek, ``On a connection of number theory with graph theory,'' \emph{Czechoslovak Math. J.} \textbf{54} (2004), 465--485; and ``Structure of digraphs associated with quadratic congruences with composite moduli,'' \emph{Discrete Math.} \textbf{306} (2006).

\bibitem{chou-shparlinski} W.-S. Chou and I.~E. Shparlinski, ``On the cycle structure of repeated exponentiation modulo a prime,'' \emph{J. Number Theory} \textbf{107} (2004), 345--356.

\bibitem{delange} H.~Delange, ``Sur la fonction sommatoire de la fonction `somme des chiffres','' \emph{Enseign. Math.} \textbf{21} (1975), 31--47.

\bibitem{drmota-poly} M.~Drmota, ``The sum of digits function of polynomial sequences,'' TU Wien. Gaussian distribution and periodic fluctuations of $S_b(P(n))$, including $P(n) = n^k$.

\bibitem{drmota-grabner} M.~Drmota and P.~Grabner, \emph{Analysis of Digital Functions and Applications}, Encyclopedia of Mathematics and Its Applications, Cambridge University Press, 2010.

\bibitem{devaney} R.~L. Devaney, \emph{An Introduction to Chaotic Dynamical Systems}, 2nd ed., Westview Press, 2003.

\bibitem{porges} A.~Porges, ``A set of eight numbers,'' \emph{American Mathematical Monthly} \textbf{52} (1945), 379--382.

\end{thebibliography}

\end{document}
